
\section{Fusion rule examples for W-algebras}
\label{sec:fusion-examples}

\subsection{\texorpdfstring{Example: Virasoro projection of $W_3(4,3)$ --- Ising model $\mathcal{M}(4,3)$}{Example: Virasoro projection of W3(4,3)}}

\begin{table}[ht]
\centering
\caption{Virasoro fusion table for $\mathcal{M}(4,3)$ (Ising model, $c = 1/2$).  After the identification $\Phi_{r,s} \sim \Phi_{q-r,p-s}$, the three Virasoro primaries are $\mathbb{I} = \Phi_{1,1}$ ($h=0$), $\sigma = \Phi_{1,2}$ ($h=1/16$), $\varepsilon = \Phi_{1,3}$ ($h=1/2$).}
\begin{tabular}{|c||c|c|c|}
\hline
$\times$ & $\mathbb{I}$ & $\sigma$ & $\varepsilon$ \\
\hline\hline
$\mathbb{I}$ & $\mathbb{I}$ & $\sigma$ & $\varepsilon$ \\
\hline
$\sigma$ & $\sigma$ & $\mathbb{I} + \varepsilon$ & $\sigma$ \\
\hline
$\varepsilon$ & $\varepsilon$ & $\sigma$ & $\mathbb{I}$ \\
\hline
\end{tabular}
\end{table}

\begin{remark}
The $W_3(4,3)$ model has central charge $c_{W_3} = 0$ and is degenerate.  Its Virasoro projection gives conformal weights $h = 0, 1/16, 1/2$, matching those of the Ising model $\mathcal{M}(4,3)$ at $c_{\mathrm{Vir}} = 1/2$, but the underlying representation categories differ.  The fusion rules above are the Virasoro fusion rules; the proper $W_3$ fusion rules require $A_2$ weight labeling (Fateev--Lukyanov).
\end{remark}

\subsection{\texorpdfstring{Example: minimal model $(6,5)$ selected rules}{Example: minimal model (6,5) selected rules}}

Key fusion products for $W_3(6,5)$:
\begin{align}
\Phi_{1,2} \times \Phi_{1,2} &= \mathbb{I} + \Phi_{2,2} + \Phi_{1,4} \\
\Phi_{2,1} \times \Phi_{2,1} &= \mathbb{I} + \Phi_{2,2} + \Phi_{4,1} \\
\Phi_{1,3} \times \Phi_{1,3} &= \mathbb{I} + \Phi_{2,2} + \Phi_{2,4} + \Phi_{4,1}
\end{align}
(In the last rule, the na\"ive Verlinde formula produces $\Phi_{1,6}$, which lies outside the fundamental domain; the $W_3$ Weyl group identification maps this to $\Phi_{4,1}$.)

\subsection{Connection to representation theory}

The fusion rules encode the tensor product structure of $W_3$ representations:
$$[\mathcal{L}_i] \otimes [\mathcal{L}_j] = \bigoplus_k N_{ij}^k [\mathcal{L}_k]$$

in the Grothendieck ring $K_0(W_3\text{-mod})$.

